%% SYNTHETIC TEST UNIT generated by writing/tools/make_outline_book.py from FINAL_OUTLINE_v2 -- not book text
\chapter{小 $x$ 与高能因子化}\label{ch:smallx}
本单元为测试用的合成文本（由 FINAL\_OUTLINE\_v2 自动生成），陈述内容从略。

\section{第 1 节（测试）}\label{sec:smallx-regge}

\begin{definition}[高能极限]\label{def:smallx-regge}
（测试用合成陈述：内容从略。）
\uses{def:prelim-wilson, def:dis-kinematics}
\end{definition}

\begin{theorem}[BFKL 方程 \tagA]\label{thm:smallx-bfkl}
（测试用合成陈述：内容从略。）
\uses{def:smallx-regge}
\end{theorem}
\begin{proofsketch}测试。\end{proofsketch}
\paragraph{地位与缺口}测试：见文献。

\section{第 2 节（测试）}\label{sec:smallx-bfkl}

\begin{proposition}[$k_T$（高能）因子化 \tagP]\label{prop:smallx-kt}
（测试用合成陈述：内容从略。）
\uses{thm:smallx-bfkl}
\end{proposition}
\paragraph{地位与缺口}测试：见文献。

\begin{theorem}[BK 方程与 JIMWLK \tagA]\label{thm:smallx-bk}
（测试用合成陈述：内容从略。）
\uses{def:smallx-regge, thm:smallx-bfkl}
\end{theorem}
\begin{proofsketch}测试。\end{proofsketch}
\paragraph{地位与缺口}测试：见文献。

\section{第 3 节（测试）}\label{sec:smallx-kt}

\begin{conjecture}[饱和与色玻璃凝聚 \tagC]\label{conj:smallx-saturation}
（测试用合成陈述：内容从略。）
\uses{thm:smallx-bk}
\end{conjecture}
\paragraph{地位与缺口}测试：见文献。

\begin{remark}[共线因子化在小 $x$ 的非一致性]\label{rem:smallx-nonuniform}
（测试用合成陈述：内容从略。）
% uses: none
\end{remark}

\section{第 4 节（测试）}\label{sec:smallx-dipole}

\begin{example}[DGLAP 与 BFKL 的双领头对数重叠]\label{ex:smallx-dla}
（测试用合成陈述：内容从略。）
\uses{prop:pdf-lo-kernels, thm:smallx-bfkl}
\end{example}
